\section{Solved Problems in Advanced Hypothesis Testing}

\subsection{Problem Statements}

\subsubsection*{1. Fundamental Level (Simple / Direct)}

\begin{problema}
	\label{prob:doc_dos_medias_iguales}
	A telemedicine company evaluates whether the mean response time for general practice consultations differs between two digital platforms (Platform 1 vs. Platform 2). Two independent random samples of times (in minutes) are obtained:
	\begin{itemize}
		\item \textbf{Platform 1:} $n_1 = 16$, $\bar{X}_1 = 14.2$ min, $S_1 = 2.4$ min.
		\item \textbf{Platform 2:} $n_2 = 18$, $\bar{X}_2 = 16.8$ min, $S_2 = 2.6$ min.
	\end{itemize}
	Assuming the populations are normal and the variances are equal ($\sigma_1^2 = \sigma_2^2$), perform a hypothesis test at the $5\%$ significance level ($\alpha = 0.05$) to determine if there is a true difference in mean response times.
\end{problema}

\begin{problema}
	\label{prob:doc_pareadas}
	In a study on the physical performance of athletes undergoing high-altitude training, maximum oxygen uptake ($\text{VO}_2$ max in ml/kg/min) was measured in a sample of $n = 12$ athletes \emph{before} ($X_1$) and \emph{after} ($X_2$) an 8-week program. The individual differences $D_i = X_{2i} - X_{1i}$ yielded a sample mean increase $\bar{D} = 3.8$ ml/kg/min and a standard deviation of differences $S_D = 2.1$ ml/kg/min.
	
	Do the data provide sufficient evidence at the $\alpha = 0.01$ level to claim that the altitude training program significantly \textbf{increases} maximum oxygen uptake?
\end{problema}

\begin{problema}
	\label{prob:doc_dos_proporciones}
	A cybersecurity company compares the effectiveness of a classical anti-spam filter against a new artificial intelligence model based on transformers. By processing two independent samples of verified incoming emails:
	\begin{itemize}
		\item \textbf{Classical Filter:} Successfully detected $X_1 = 340$ out of $n_1 = 400$ malicious emails as spam.
		\item \textbf{AI Model:} Successfully detected $X_2 = 465$ out of $n_2 = 500$ malicious emails as spam.
	\end{itemize}
	Is there evidence at the $\alpha = 0.02$ level to state that the AI Model has a \textbf{significantly higher detection rate} than the Classical Filter?
\end{problema}

\subsubsection*{2. Operational Level (Intermediate / Developmental)}

\begin{problema}
	\label{prob:doc_dos_varianzas}
	Before performing a pooled $t$-test comparing the means of two analog sensors, one must verify whether the homoscedasticity assumption ($H_0: \sigma_1^2 = \sigma_2^2$) holds. Two independent random samples are taken, yielding sample variances $S_1^2 = 14.8$ ($n_1 = 16$) and $S_2^2 = 5.2$ ($n_2 = 11$).
	
	Perform a two-tailed $F$-test at the $\alpha = 0.10$ level to test the hypothesis of equal variances. Is it justifiable to assume homoscedasticity?
\end{problema}

\begin{problema}
	\label{prob:homogeneidad_contingencia}
	A market research firm analyzes the preferences of three generational cohorts of software users (Generation Z, Millennials, and Generation X) regarding their preferred monetization model in digital tools: Monthly subscription, One-time payment (perpetual license), or Freemium with ads. Three independent samples of fixed sizes $n_1 = 200$, $n_2 = 250$, and $n_3 = 150$ are selected, obtaining the observed frequencies shown in Table \ref{tab:homog_obs}.
	\begin{table}[h!]
		\centering
		\begin{tabular}{lcccc}
			\hline
			\textbf{Preferred Model} & \textbf{Gen Z ($j=1$)} & \textbf{Millennials ($j=2$)} & \textbf{Gen X ($j=3$)} & \textbf{Row Total ($R_i$)} \\ \hline
			Subscription ($i=1$) & 80 & 120 & 45 & \textbf{245} \\
			One-time Payment ($i=2$) & 30 & 65 & 70 & \textbf{165} \\
			Freemium ($i=3$) & 90 & 65 & 35 & \textbf{190} \\ \hline
			\textbf{Column Total ($C_j$)} & \textbf{200} & \textbf{250} & \textbf{150} & \textbf{600} \\ \hline
		\end{tabular}
		\caption{Observed frequencies ($O_{ij}$) by monetization model across three generational cohorts.}
		\label{tab:homog_obs}
	\end{table}
	
	Perform a test of homogeneity of proportions at the $\alpha = 0.01$ level to determine whether the distribution of monetization preferences is identical across the three generations.
\end{problema}

\begin{problema}
	\label{prob:doc_welch_hetero}
	A data scientist compares the intermediate execution times (in seconds) consumed by two sparse clustering algorithms. Since the hardware architecture of the second algorithm generates far more volatility, the population variances are explicitly heterogeneous ($\sigma_1^2 \neq \sigma_2^2$). The following experimental measurements are extracted:
	\begin{itemize}
		\item \textbf{Algorithm 1:} $n_1 = 15$, $\bar{X}_1 = 85.2$ s, $S_1 = 12.4$ s.
		\item \textbf{Algorithm 2:} $n_2 = 18$, $\bar{X}_2 = 74.5$ s, $S_2 = 5.2$ s.
	\end{itemize}
	\begin{enumerate}
		\item Calculate the effective degrees of freedom ($\nu$) using the Welch-Satterthwaite equation.
		\item Execute the two-tailed Welch $t$-test with $\alpha = 0.05$ to determine if the mean execution times of both algorithms differ statistically significantly.
	\end{enumerate}
\end{problema}

\subsubsection*{3. Analytical Level (Advanced / Connection)}

\begin{problema}
	\label{prob:doc_k_proporciones_marascuilo}
	To evaluate the level of adherence to ethical standards on algorithmic bias, a regulatory audit agency examines samples of machine learning models deployed by $k = 4$ corporate technology firms (Companies A, B, C, and D). It is verified whether each model satisfies fairness criteria (Success = Complies, Failure = Does not comply):
	\begin{itemize}
		\item \textbf{Company A:} $X_1 = 82$ comply, out of $n_1 = 100$ models ($\hat{p}_1 = 0.820$).
		\item \textbf{Company B:} $X_2 = 64$ comply, out of $n_2 = 100$ models ($\hat{p}_2 = 0.640$).
		\item \textbf{Company C:} $X_3 = 90$ comply, out of $n_3 = 120$ models ($\hat{p}_3 = 0.750$).
		\item \textbf{Company D:} $X_4 = 44$ comply, out of $n_4 = 80$ models ($\hat{p}_4 = 0.550$).
	\end{itemize}
	\begin{enumerate}
		\item Perform an omnibus $\chi^2$ test for equal proportions ($H_0: p_1=p_2=p_3=p_4$) at the $\alpha = 0.05$ level.
		\item If the omnibus null hypothesis is rejected, apply the \textbf{Marascuilo post-hoc multiple comparison procedure} to determine which pairs of companies differ significantly from one another at family-wise level $\alpha = 0.05$.
	\end{enumerate}
\end{problema}

\begin{problema}
	\label{prob:doc_errores_potencia_muestra}
	\textbf{Comprehensive analysis of Type I Error ($\alpha$), Type II Error ($\beta$), and Statistical Power ($1-\beta$):}
	The server response time for transaction processing at a fintech follows a normal distribution with known population standard deviation $\sigma = 15$ milliseconds. The operating standard requires that mean latency not exceed $100$ ms ($H_0: \mu \le 100$ vs. $H_a: \mu > 100$). A random daily sample of $n = 25$ transactions is monitored at a significance level of $\alpha = 0.05$.
	\begin{enumerate}
		\item Determine the decision rule on the original scale of the sample mean latency ($\bar{X}_c = \mu_0 + Z_\alpha \frac{\sigma}{\sqrt{n}}$).
		\item Suppose that due to network congestion, the true mean latency of the population silently increased to $\mu_a = 108$ ms. Calculate the exact probability of committing a \textbf{Type II Error ($\beta$)}—failing to detect the congestion—and determine the exact value of the \textbf{Power of the test ($1-\beta$)}.
		\item If the VP of Engineering demands that the test maintain risk $\alpha = 0.05$ while being capable of detecting a true congestion of $\mu_a = 108$ ms with a \textbf{Statistical Power of $90\%$} ($1-\beta = 0.90 \implies Z_\beta = 1.282$), derive and precisely calculate what sample size $n$ the system must process in each audit.
	\end{enumerate}
\end{problema}

\subsubsection*{4. Challenging Level (Experts / Deep Deduction)}

\begin{problema}
	\label{prob:doc_desaf_lrt_wilks}
	\textbf{Derivation of the Likelihood Ratio Test (LRT) for variances and invocation of Wilks' Theorem:}
	Let two independent random samples be $X_1, \dots, X_{n_1} \sim N(\mu_1, \sigma_1^2)$ and $Y_1, \dots, Y_{n_2} \sim N(\mu_2, \sigma_2^2)$ where both means and variances are completely unknown. We wish to test the parametric hypothesis of homoscedasticity $H_0: \sigma_1^2 = \sigma_2^2 = \sigma^2$ against $H_a: \sigma_1^2 \neq \sigma_2^2$.
	\begin{enumerate}
		\item Derive the logarithm of the joint bivariate likelihood function $\ln L(\mu_1, \mu_2, \sigma_1^2, \sigma_2^2)$ under the unrestricted total parameter space $\Theta$, explicitly writing out the maximum likelihood dispersion estimators $\hat{\sigma}_1^2 = \frac{1}{n_1}\sum (X_i - \bar{X})^2$ and $\hat{\sigma}_2^2 = \frac{1}{n_2}\sum (Y_i - \bar{Y})^2$.
		\item Derive the log-likelihood under the restricted null subspace $\Theta_0$ (where $\sigma_1^2 = \sigma_2^2 = \sigma_0^2$), demonstrating that the combined maximum likelihood estimate is:
		\begin{align}
			\hat{\sigma}_0^2 = \frac{n_1 \hat{\sigma}_1^2 + n_2 \hat{\sigma}_2^2}{n_1 + n_2}.
		\end{align}
		\item Prove that the logarithmic likelihood ratio statistic $\Lambda = \dfrac{L(\Theta_0)}{L(\Theta)}$ satisfies the exact algebraic identity:
		\begin{align}
			-2 \ln \Lambda = n_1 \ln\left( \frac{\hat{\sigma}_0^2}{\hat{\sigma}_1^2} \right) + n_2 \ln\left( \frac{\hat{\sigma}_0^2}{\hat{\sigma}_2^2} \right).
		\end{align}
		\item Explain why, under the null hypothesis, the sampling distribution of $-2 \ln \Lambda$ converges asymptotically to $\chi^2_1$ according to \textbf{Samuel S. Wilks' Theorem}, justifying the computation of lost parametric dimension between $\Theta$ and $\Theta_0$.
	\end{enumerate}
\end{problema}

\begin{problema}
	\label{prob:doc_desaf_pvalor_teorema}
	\textbf{Probabilistic structure, stochastic dominance, and analytical geometry of the $p$-value:}
	In the data science literature, it is common to confuse the observed $p$-value with the posterior probability of the null hypothesis ($P(H_0 \mid \text{data})$). Formalize the true nature of the $p$-value by treating it as a random variable $P(X)$ over the space of observed samples:
	\begin{enumerate}
		\item Consider a continuous univariate test whose test statistic $T(X)$ has a strictly monotonic continuous distribution function $F_0(t) = P_{H_0}(T \le t)$ and survival function $S_0(t) = 1 - F_0(t) = P_{H_0}(T \ge t)$ under the null hypothesis. The random $p$-value for upper-tailed tests is defined by $P = S_0(T)$.
		Rigorously prove using differentiation and the Probability Integral Transform Theorem that, \textbf{when $H_0$ is true, the random variable $p$-value $P$ is distributed exactly Uniformly over the unit interval $[0, 1]$} ($P \sim U(0,1)$), satisfying $P_{H_0}(P \le \alpha) = \alpha$ for all $\alpha \in (0,1)$.
		\item When the alternative hypothesis $H_a$ is true (a true effect size exists in nature), the statistic adopts a different alternative distribution $F_a(t) = P_{H_a}(T \le t)$. Algebraically prove that the cumulative distribution function of the $p$-value under $H_a$, denoted by $G_P(u) = P_{H_a}(P \le u)$ for $u \in (0,1)$, is a \textbf{strictly concave} function that stochastically dominates the uniform diagonal ($G_P(u) > u$).
		\item Prove why the differential density of the $p$-value at the boundary of absolute rejection $\lim_{u \to 0^+} g_P(u) = \lim_{u \to 0^+} G_P'(u)$ is directly proportional to the \textbf{monotone likelihood ratio} $\dfrac{f_a(t)}{f_0(t)}$ in the extreme region $t \to \infty$, justifying why, in the presence of any real effect, experimental $p$-values massively concentrate in the immediate neighborhood of zero.
	\end{enumerate}
\end{problema}


\subsection{Detailed Solutions}

\subsubsection*{1. Fundamental Level (Simple / Direct)}

\begin{solucion}
	\textbf{Step 1: Hypothesis formulation.}
	We want to test the equality of both means against a two-tailed difference:
	\begin{align}
		H_0&: \mu_1 - \mu_2 = 0 \\
		H_a&: \mu_1 - \mu_2 \neq 0
	\end{align}
	
	\textbf{Step 2: Calculation of the pooled variance ($S_p^2$).}
	The degrees of freedom are $\nu = n_1 + n_2 - 2 = 16 + 18 - 2 = 32$.
	\begin{align}
		S_p^2 &= \frac{(n_1 - 1)S_1^2 + (n_2 - 1)S_2^2}{n_1 + n_2 - 2} = \frac{(15)(2.4)^2 + (17)(2.6)^2}{32} \\
		&= \frac{15(5.76) + 17(6.76)}{32} = \frac{86.4 + 114.92}{32} = \frac{201.32}{32} \approx 6.29125 \text{ min}^2.
	\end{align}
	The pooled standard deviation is $S_p = \sqrt{6.29125} \approx 2.508$ min.
	
	\textbf{Step 3: Calculation of the test statistic $t$.}
	\begin{align}
		t = \frac{(\bar{X}_1 - \bar{X}_2) - 0}{S_p \sqrt{\frac{1}{n_1} + \frac{1}{n_2}}} = \frac{14.2 - 16.8}{2.508 \sqrt{\frac{1}{16} + \frac{1}{18}}} = \frac{-2.6}{2.508 \sqrt{0.0625 + 0.05556}} = \frac{-2.6}{2.508 \times 0.3436} \approx \frac{-2.6}{0.8617} \approx -3.017.
	\end{align}
	
	\textbf{Step 4: Decision rule and conclusion.}
	For a two-tailed test with $\alpha = 0.05$ and $\nu = 32$ degrees of freedom, the critical value in the $t$-table is $t_{0.025, \, 32} \approx 2.037$.
	The rejection region consists of values $t < -2.037$ or $t > 2.037$.
	
	Since our calculated statistic $t = -3.017$ clearly falls into the left-tailed rejection region ($-3.017 < -2.037$), \textbf{we reject $H_0$}. We conclude at the $5\%$ level that there is a statistically significant difference in mean response times between the two platforms, with Platform 1 being significantly faster.
\end{solucion}

\begin{solucion}
	\textbf{Step 1: Formulation of one-tailed hypotheses (paired).}
	Since we evaluate whether the program increases the post-training value relative to the pre-training value ($X_2 > X_1 \implies D > 0$), we formulate a right-tailed test on the population mean difference $\mu_D$:
	\begin{align}
		H_0&: \mu_D \leq 0 \\
		H_a&: \mu_D > 0
	\end{align}
	
	\textbf{Step 2: Calculation of the $t$ statistic for paired samples.}
	The degrees of freedom are $\nu = n - 1 = 12 - 1 = 11$.
	\begin{align}
		t = \frac{\bar{D} - 0}{S_D / \sqrt{n}} = \frac{3.8}{2.1 / \sqrt{12}} = \frac{3.8}{2.1 / 3.4641} = \frac{3.8}{0.6062} \approx 6.268.
	\end{align}
	
	\textbf{Step 3: Decision rule and conclusion.}
	For an upper one-tailed test with $\alpha = 0.01$ and $\nu = 11$, the critical value of Student's $t$-distribution is:
	\begin{align}
		t_{0.01, \, 11} = 2.718.
	\end{align}
	The rejection region is $t > 2.718$. Given that our observed value $t = 6.268 \gg 2.718$, \textbf{we categorically reject $H_0$}. There is overwhelming statistical evidence at the $1\%$ level to affirm that the high-altitude training program produces a real increase in athletes' maximum $\text{VO}_2$.
\end{solucion}

\begin{solucion}
	\textbf{Step 1: Formulation of hypotheses.}
	Let $p_1$ and $p_2$ denote the population detection proportions of the classical filter and the AI model, respectively. We evaluate whether $p_2 > p_1$, that is, if $p_1 - p_2 < 0$:
	\begin{align}
		H_0&: p_1 - p_2 \geq 0 \\
		H_a&: p_1 - p_2 < 0 \quad \text{(left-tailed test)}.
	\end{align}
	
	\textbf{Step 2: Empirical proportions and pooled proportion ($\bar{p}$).}
	The individual proportions are $\hat{p}_1 = 340/400 = 0.850$ ($85\%$) and $\hat{p}_2 = 465/500 = 0.930$ ($93\%$).
	Under the null hypothesis that both proportions are equal, the pooled estimator combining both samples is:
	\begin{align}
		\bar{p} = \frac{X_1 + X_2}{n_1 + n_2} = \frac{340 + 465}{400 + 500} = \frac{805}{900} \approx 0.89444.
	\end{align}
	
	\textbf{Step 3: Calculation of the test statistic $Z$.}
	The standard error under $H_0$ is:
	\begin{align}
		SE_0 = \sqrt{\bar{p}(1-\bar{p})\left(\frac{1}{n_1} + \frac{1}{n_2}\right)} &= \sqrt{0.89444(0.10556)\left(\frac{1}{400} + \frac{1}{500}\right)} \\
		&= \sqrt{0.094416 \times 0.0045} = \sqrt{0.00042487} \approx 0.02061.
	\end{align}
	The standardized normal statistic becomes:
	\begin{align}
		Z = \frac{\hat{p}_1 - \hat{p}_2}{SE_0} = \frac{0.850 - 0.930}{0.02061} = \frac{-0.080}{0.02061} \approx -3.881.
	\end{align}
	
	\textbf{Step 4: Decision and conclusion.}
	For a left-tailed test at the $2\%$ level ($\alpha = 0.02$), we locate the lower normal quantile:
	\begin{align}
		-Z_{0.02} = -2.054.
	\end{align}
	Since the observed statistic $Z = -3.881$ falls into the rejection region ($Z < -2.054$, with $p$-value $< 0.0001$), \textbf{we reject $H_0$}. We demonstrate at the $2\%$ significance level that the new AI Model has a superior detection rate compared to the Classical Filter.
\end{solucion}

\subsubsection*{2. Operational Level (Intermediate / Developmental)}

\begin{solucion}
	\textbf{Step 1: Formulation of two-tailed hypotheses.}
	\begin{align}
		H_0&: \sigma_1^2 = \sigma_2^2 \\
		H_a&: \sigma_1^2 \neq \sigma_2^2
	\end{align}
	
	\textbf{Step 2: Calculation of the $F$ statistic.}
	Placing the larger variance in the numerator to compare against the upper tail of the Snedecor table:
	\begin{align}
		F = \frac{S_1^2}{S_2^2} = \frac{14.8}{5.2} \approx 2.846.
	\end{align}
	The degrees of freedom are $\nu_1 = n_1 - 1 = 15$ in the numerator and $\nu_2 = n_2 - 1 = 10$ in the denominator.
	
	\textbf{Step 3: Obtaining the critical points of the Fisher-Snedecor distribution.}
	For a two-tailed test at level $\alpha = 0.10$, the area in each extreme tail is $\alpha/2 = 0.05$. Consulting the $F$-table for $15$ and $10$ degrees of freedom:
	\begin{align}
		F_{0.05, \, 15, \, 10} = 2.85.
	\end{align}
	The lower critical value is $F_{0.95, \, 15, \, 10} = \dfrac{1}{F_{0.05, \, 10, \, 15}} = \dfrac{1}{2.54} \approx 0.394$.
	
	\textbf{Step 4: Technical conclusion.}
	Our observed statistic $F = 2.846$ is right on the boundary of the upper rejection threshold ($2.846 \approx 2.85$). Strictly speaking ($2.846 < 2.85$), $H_0$ is not formally rejected at the $10\%$ level. However, because it lies directly on the rejection border, assuming pure homoscedasticity in subsequent steps is inadvisable, and applying Welch's heteroscedastic $t$-test is strongly recommended.
\end{solucion}

\begin{solucion}
	\textbf{Step 1: Formulation of homogeneity hypotheses.}
	\begin{align}
		H_0&: \text{The distribution of preferred model is homogeneous across all three generations } (p_{i1}=p_{i2}=p_{i3} \; \forall i). \\
		H_a&: \text{Preference distributions differ in at least one generation.}
	\end{align}
	
	\textbf{Step 2: Calculation of expected frequencies under $H_0$ ($E_{ij} = \frac{R_i C_j}{N}$).}
	The grand total is $N = 600$. We calculate cell by cell:
	\begin{itemize}
		\item $E_{11} = \frac{245 \times 200}{600} = 81.67$, \quad $E_{12} = \frac{245 \times 250}{600} = 102.08$, \quad $E_{13} = \frac{245 \times 150}{600} = 61.25$.
		\item $E_{21} = \frac{165 \times 200}{600} = 55.00$, \quad $E_{22} = \frac{165 \times 250}{600} = 68.75$, \quad $E_{23} = \frac{165 \times 150}{600} = 41.25$.
		\item $E_{31} = \frac{190 \times 200}{600} = 63.33$, \quad $E_{32} = \frac{190 \times 250}{600} = 79.17$, \quad $E_{33} = \frac{190 \times 150}{600} = 47.50$.
	\end{itemize}
	We verify that all $E_{ij} \geq 5$, confirming that the asymptotic $\chi^2$ approximation is valid.
	
	\textbf{Step 3: Calculation of the $\chi^2$ test statistic.}
	We calculate the contribution $\dfrac{(O_{ij}-E_{ij})^2}{E_{ij}}$ for each of the 9 cells:
	\begin{align}
		\chi^2 &= \frac{(80-81.67)^2}{81.67} + \frac{(120-102.08)^2}{102.08} + \frac{(45-61.25)^2}{61.25} + \frac{(30-55)^2}{55.00} + \frac{(65-68.75)^2}{68.75} + \frac{(70-41.25)^2}{41.25} \nonumber \\
		&\quad + \frac{(90-63.33)^2}{63.33} + \frac{(65-79.17)^2}{79.17} + \frac{(35-47.50)^2}{47.50} \\
		&= 0.034 + 3.146 + 4.311 + 11.364 + 0.205 + 20.030 + 11.231 + 2.536 + 3.290 = 56.147.
	\end{align}
	
	\textbf{Step 4: Degrees of freedom and decision.}
	For a $3 \times 3$ table, $\nu = (3-1)(3-1) = 4$. The critical value at $\alpha = 0.01$ is $\chi^2_{0.01, \, 4} = 13.277$.
	Since $\chi^2 = 56.147 \gg 13.277$ ($p$-value $< 0.00001$), \textbf{we reject $H_0$}. The distribution of monetization preferences is not identical; Generation X strongly prefers one-time payment ($70$ vs. $41.25$ expected) while Gen Z overwhelmingly favors the Freemium model ($90$ vs. $63.33$ expected).
\end{solucion}

\begin{solucion}
	\begin{enumerate}
		\item \textbf{Calculation of Welch degrees of freedom ($\nu$):}
		The squared sample errors of the mean are:
		\begin{align}
			\frac{S_1^2}{n_1} = \frac{12.4^2}{15} = \frac{153.76}{15} \approx 10.2507, \quad \frac{S_2^2}{n_2} = \frac{5.2^2}{18} = \frac{27.04}{18} \approx 1.5022.
		\end{align}
		Applying the analytical Welch-Satterthwaite approximation:
		\begin{align}
			\nu = \frac{\left( \dfrac{S_1^2}{n_1} + \dfrac{S_2^2}{n_2} \right)^2}{\dfrac{(S_1^2 / n_1)^2}{n_1 - 1} + \dfrac{(S_2^2 / n_2)^2}{n_2 - 1}} = \frac{(10.2507 + 1.5022)^2}{\dfrac{10.2507^2}{14} + \dfrac{1.5022^2}{17}} = \frac{(11.7529)^2}{\dfrac{105.0769}{14} + \dfrac{2.2566}{17}} = \frac{138.1307}{7.5055 + 0.1327} \approx 18.08.
		\end{align}
		Taking the conservative integer floor, we set $\nu = 18$ degrees of freedom.
		
		\item \textbf{Calculation of the two-tailed $t$ statistic at $\alpha = 0.05$:}
		The standard error of the difference for heterogeneous variances is:
		\begin{align}
			\text{SE}_W = \sqrt{ \frac{S_1^2}{n_1} + \frac{S_2^2}{n_2} } = \sqrt{ 10.2507 + 1.5022 } = \sqrt{11.7529} \approx 3.4282 \text{ s.}
		\end{align}
		The elementary difference between observed sample means is $\bar{X}_1 - \bar{X}_2 = 85.2 - 74.5 = 10.7$ s. The standardized $t$ statistic becomes:
		\begin{align}
			t = \frac{\bar{X}_1 - \bar{X}_2}{\text{SE}_W} = \frac{10.7}{3.4282} \approx 3.1212.
		\end{align}
		For a two-tailed test with $\alpha = 0.05$ and $\nu = 18$, Student's quantile is $t_{0.025, \, 18} = 2.101$.
		Because $t = 3.121 > 2.101$, \textbf{we reject $H_0$}. We conclude at 95\% confidence that Algorithm 2 consumes a significantly lower mean execution time than Algorithm 1.
	\end{enumerate}
\end{solucion}

\subsubsection*{3. Analytical Level (Advanced / Connection)}

\begin{solucion}
	\textbf{Part 1: Omnibus $\chi^2$ test for equality across 4 proportions.}
	The total number of successes across the $k=4$ samples is $\sum X_j = 82 + 64 + 90 + 44 = 280$. Total observations equal $\sum n_j = 100 + 100 + 120 + 80 = 400$. The pooled overall proportion is:
	\begin{align}
		\bar{p} = \frac{280}{400} = 0.70 \quad (70\%).
	\end{align}
	We apply the direct formula for the multiple proportion $\chi^2$ statistic:
	\begin{align}
		\chi^2 = \sum_{j=1}^4 \frac{(X_j - n_j \bar{p})^2}{n_j \bar{p}(1-\bar{p})}.
	\end{align}
	The denominator scales with sample weight: $n_j \bar{p}(1-\bar{p}) = n_j(0.70)(0.30) = 0.21 n_j$. We evaluate cell by cell:
	\begin{itemize}
		\item \textbf{Company A:} $n_1 \bar{p} = 70$. Contribution: $\frac{(82 - 70)^2}{21} = \frac{144}{21} \approx 6.8571$.
		\item \textbf{Company B:} $n_2 \bar{p} = 70$. Contribution: $\frac{(64 - 70)^2}{21} = \frac{36}{21} \approx 1.7143$.
		\item \textbf{Company C:} $n_3 \bar{p} = 84$. Contribution: $\frac{(90 - 84)^2}{120 \times 0.21} = \frac{36}{25.2} \approx 1.4286$.
		\item \textbf{Company D:} $n_4 \bar{p} = 56$. Contribution: $\frac{(44 - 56)^2}{80 \times 0.21} = \frac{144}{16.8} \approx 8.5714$.
	\end{itemize}
	Summing all 4 contributions:
	\begin{align}
		\chi^2 = 6.8571 + 1.7143 + 1.4286 + 8.5714 = 18.5714.
	\end{align}
	The degrees of freedom are $\nu = k - 1 = 4 - 1 = 3$. For $\alpha = 0.05$, the quantile is $\chi^2_{0.05, \, 3} = 7.815$.
	Since $\chi^2 = 18.5714 > 7.815$, \textbf{we reject $H_0$}. Ethical compliance rates differ significantly among the four companies.
	
	\textbf{Part 2: Marascuilo post-hoc procedure.}
	The Marascuilo critical value uses the square root of the omnibus statistic: $\sqrt{\chi^2_{0.05, \, 3}} = \sqrt{7.815} \approx 2.7955$.
	For each pair $(i, j)$, we evaluate the difference $|\hat{p}_i - \hat{p}_j|$ against the pairwise critical range $r_{ij} = 2.7955 \sqrt{\frac{\hat{p}_i(1-\hat{p}_i)}{n_i} + \frac{\hat{p}_j(1-\hat{p}_j)}{n_j}}$:
	\begin{itemize}
		\item \textbf{A vs. B:} $|\hat{p}_1 - \hat{p}_2| = |0.82 - 0.64| = 0.180$.
		$r_{12} = 2.7955 \sqrt{\frac{0.82(0.18)}{100} + \frac{0.64(0.36)}{100}} = 2.7955 \sqrt{0.00378} \approx 0.1719$.
		Since $0.180 > 0.1719$, \textbf{the A vs. B difference is statistically significant}.
		\item \textbf{A vs. C:} $|\hat{p}_1 - \hat{p}_3| = |0.82 - 0.75| = 0.070$.
		$r_{13} = 2.7955 \sqrt{0.001476 + \frac{0.75(0.25)}{120}} \approx 0.1541$. ($0.070 \le 0.1541 \implies \text{Not significant}$).
		\item \textbf{A vs. D:} $|\hat{p}_1 - \hat{p}_4| = |0.82 - 0.55| = 0.270$.
		$r_{14} = 2.7955 \sqrt{0.001476 + \frac{0.55(0.45)}{80}} \approx 0.1890$. ($0.270 > 0.1890 \implies \textbf{Significant}$).
		\item \textbf{B vs. C:} $|\hat{p}_2 - \hat{p}_3| = |0.64 - 0.75| = 0.110$.
		$r_{23} \approx 0.1738$. ($0.110 \le 0.1738 \implies \text{Not significant}$).
		\item \textbf{B vs. D:} $|\hat{p}_2 - \hat{p}_4| = |0.64 - 0.55| = 0.090$.
		$r_{24} \approx 0.2054$. ($0.090 \le 0.2054 \implies \text{Not significant}$).
		\item \textbf{C vs. D:} $|\hat{p}_3 - \hat{p}_4| = |0.75 - 0.55| = 0.200$.
		$r_{34} = 2.7955 \sqrt{\frac{0.75(0.25)}{120} + \frac{0.55(0.45)}{80}} = 2.7955 \sqrt{0.004656} \approx 0.1908$. ($0.200 > 0.1908 \implies \textbf{Significant}$).
	\end{itemize}
	\textbf{Conclusion:} The pairs showing proven real differences are \textbf{A vs. B} ($82\%$ vs. $64\%$), \textbf{A vs. D} ($82\%$ vs. $55\%$), and \textbf{C vs. D} ($75\%$ vs. $55\%$).
\end{solucion}

\begin{solucion}
	\begin{enumerate}
		\item \textbf{Decision rule and critical boundary on the original scale ($\bar{X}_c$):}
		The standard error of the population mean is $\text{SE}(\bar{X}) = \dfrac{\sigma}{\sqrt{n}} = \dfrac{15}{\sqrt{25}} = \dfrac{15}{5} = 3$ ms.
		For an upper one-tailed test at significance level $\alpha = 0.05$, the critical normal quantile is $Z_{0.05} = 1.645$.
		The critical decision boundary where $H_0$ is rejected is obtained by solving:
		\begin{align}
			\bar{X}_c = \mu_0 + Z_{0.05} \cdot \text{SE}(\bar{X}) = 100 + 1.645(3) = 100 + 4.935 = 104.935 \text{ ms.}
		\end{align}
		\textbf{Operational rule:} If in any daily audit the sample mean transaction latency exceeds $104.935$ ms ($\bar{X} > 104.935$), the system rejects $H_0$ and triggers a network congestion alert.
		
		\item \textbf{Calculation of Type II Error probability ($\beta$) and Power ($1-\beta$) when $\mu_a = 108$ ms:}
		A \textbf{Type II Error ($\beta$)} occurs when we erroneously accept $H_0$ ($\bar{X} \le \bar{X}_c$) while in the real world the alternative hypothesis is true ($\mu = 108$ ms):
		\begin{align}
			\beta = P\left(\bar{X} \le 104.935 \mid \mu = 108\right).
		\end{align}
		We standardize the critical threshold relative to the true alternative distribution centered at $108$:
		\begin{align}
			Z_\beta = \frac{\bar{X}_c - \mu_a}{\sigma/\sqrt{n}} = \frac{104.935 - 108}{3} = \frac{-3.065}{3} \approx -1.0217.
		\end{align}
		We evaluate the cumulative left-hand probability in standard normal tables:
		\begin{align}
			\beta = \Phi(-1.0217) = 1 - \Phi(1.0217) \approx 1 - 0.8465 = 0.1535 \text{ (15.35\%).}
		\end{align}
		The \textbf{Statistical Power ($1-\beta$)}—the real probability of successfully detecting the congestion—is obtained by direct subtraction:
		\begin{align}
			\text{Power} = 1 - \beta = 1 - 0.1535 = 0.8465 \text{ (84.65\%).}
		\end{align}
		
		\item \textbf{Derivation of the minimum sample size $n$ for a target power of $90\%$ ($1-\beta = 0.90$):}
		To guarantee that the rejection boundary geometrically separates the upper area $\alpha = 0.05$ of the null curve centered at $\mu_0 = 100$ and the left area $\beta = 0.10$ of the alternative curve centered at $\mu_a = 108$ ($Z_\beta = 1.282$), the following equality of standardized distances is required:
		\begin{align}
			\bar{X}_c = \mu_0 + Z_\alpha \frac{\sigma}{\sqrt{n}} = \mu_a - Z_\beta \frac{\sigma}{\sqrt{n}}.
		\end{align}
		Combining standard error terms on one side:
		\begin{align}
			(Z_\alpha + Z_\beta)\frac{\sigma}{\sqrt{n}} = \mu_a - \mu_0 \implies \sqrt{n} = \frac{(Z_\alpha + Z_\beta)\sigma}{\mu_a - \mu_0}.
		\end{align}
		Squaring both sides yields the analytical sample size formula for simultaneous control of $\alpha$ and $\beta$ errors:
		\begin{align}
			n = \left( \frac{(Z_\alpha + Z_\beta)\sigma}{\mu_a - \mu_0} \right)^2.
		\end{align}
		Substituting the fintech's numerical parameters ($Z_{0.05} = 1.645, Z_{0.10} = 1.282, \sigma = 15, \Delta = 108 - 100 = 8$):
		\begin{align}
			n = \left( \frac{(1.645 + 1.282) \times 15}{8} \right)^2 = \left( \frac{2.927 \times 15}{8} \right)^2 = \left( \frac{43.905}{8} \right)^2 = (5.488125)^2 \approx 30.1195.
		\end{align}
		Applying conservative ceiling rounding:
		\begin{align}
			n_{\text{optimal}} = \lceil 30.12 \rceil = 31 \text{ daily transactions.}
		\end{align}
	\end{enumerate}
\end{solucion}

\subsubsection*{4. Challenging Level (Experts / Deep Deduction)}

\begin{solucion}
	\begin{enumerate}
		\item \textbf{Unrestricted log-likelihood ($\Theta$) and maximum likelihood estimators (MLE):}
		Given two independent normal samples of sizes $n_1$ and $n_2$, the joint density or likelihood function $L(\Theta) = L(\mu_1, \mu_2, \sigma_1^2, \sigma_2^2)$ is the product of individual densities:
		\begin{align}
			L(\Theta) = \left[ (2\pi\sigma_1^2)^{-n_1/2} \exp\left(-\frac{1}{2\sigma_1^2}\sum_{i=1}^{n_1}(X_i - \mu_1)^2\right) \right] \times \left[ (2\pi\sigma_2^2)^{-n_2/2} \exp\left(-\frac{1}{2\sigma_2^2}\sum_{j=1}^{n_2}(Y_j - \mu_2)^2\right) \right].
		\end{align}
		Taking the natural logarithm of the entire likelihood expression:
		\begin{align}
			\ln L(\Theta) = -\frac{n_1}{2}\ln(2\pi) - \frac{n_1}{2}\ln(\sigma_1^2) - \frac{1}{2\sigma_1^2}\sum (X_i - \mu_1)^2 -\frac{n_2}{2}\ln(2\pi) - \frac{n_2}{2}\ln(\sigma_2^2) - \frac{1}{2\sigma_2^2}\sum (Y_j - \mu_2)^2.
		\end{align}
		Differentiating with respect to each parameter and setting equal to zero, the maximum likelihood estimators across the full space $\Theta$ are: $\hat{\mu}_1 = \bar{X}, \hat{\mu}_2 = \bar{Y}$, and the maximum likelihood variances:
		\begin{align}
			\hat{\sigma}_1^2 = \frac{1}{n_1}\sum (X_i - \bar{X})^2, \quad \hat{\sigma}_2^2 = \frac{1}{n_2}\sum (Y_j - \bar{Y})^2.
		\end{align}
		Substituting the identities $\sum(X_i - \bar{X})^2 = n_1\hat{\sigma}_1^2$ and $\sum(Y_j - \bar{Y})^2 = n_2\hat{\sigma}_2^2$ back into $\ln L(\Theta)$, the exponential terms simplify algebraically to exact scalars ($\frac{n_1\hat{\sigma}_1^2}{2\hat{\sigma}_1^2} = \frac{n_1}{2}$):
		\begin{align}
			\ln L(\hat{\Theta}) = -\frac{n_1 + n_2}{2}\ln(2\pi) - \frac{n_1}{2}\ln(\hat{\sigma}_1^2) - \frac{n_2}{2}\ln(\hat{\sigma}_2^2) - \frac{n_1 + n_2}{2}.
		\end{align}
		
		\item \textbf{Log-likelihood under the null subspace ($\Theta_0: \sigma_1^2 = \sigma_2^2 = \sigma_0^2$):}
		Under the homoscedasticity null hypothesis, variance is a single common parameter $\sigma_0^2$. The log-likelihood is:
		\begin{align}
			\ln L(\Theta_0) = -\frac{n_1 + n_2}{2}\ln(2\pi) - \frac{n_1 + n_2}{2}\ln(\sigma_0^2) - \frac{1}{2\sigma_0^2} \left[ \sum(X_i - \mu_1)^2 + \sum(Y_j - \mu_2)^2 \right].
		\end{align}
		Maximizing with respect to $\sigma_0^2$ (with means at their sample optima $\bar{X}, \bar{Y}$), we obtain the pooled likelihood variance:
		\begin{align}
			\hat{\sigma}_0^2 = \frac{\sum(X_i - \bar{X})^2 + \sum(Y_j - \bar{Y})^2}{n_1 + n_2} = \frac{n_1 \hat{\sigma}_1^2 + n_2 \hat{\sigma}_2^2}{n_1 + n_2}.
		\end{align}
		Evaluating the null likelihood at this supremum by combining exponents:
		\begin{align}
			\ln L(\hat{\Theta}_0) = -\frac{n_1 + n_2}{2}\ln(2\pi) - \frac{n_1 + n_2}{2}\ln(\hat{\sigma}_0^2) - \frac{n_1 + n_2}{2}.
		\end{align}
		
		\item \textbf{Derivation of the logarithmic likelihood ratio statistic ($-2\ln \Lambda$):}
		The likelihood ratio is defined as the ratio of likelihood suprema across both spaces: $\Lambda = \dfrac{L(\hat{\Theta}_0)}{L(\hat{\Theta})}$. Taking logarithms and subtracting $\ln L(\hat{\Theta}_0) - \ln L(\hat{\Theta})$:
		\begin{align}
			\ln \Lambda &= \left[ -\frac{n_1+n_2}{2}\ln(2\pi) - \frac{n_1+n_2}{2}\ln(\hat{\sigma}_0^2) - \frac{n_1+n_2}{2} \right] - \left[ -\frac{n_1+n_2}{2}\ln(2\pi) - \frac{n_1}{2}\ln(\hat{\sigma}_1^2) - \frac{n_2}{2}\ln(\hat{\sigma}_2^2) - \frac{n_1+n_2}{2} \right] \nonumber \\
			&= -\frac{n_1 + n_2}{2}\ln(\hat{\sigma}_0^2) + \frac{n_1}{2}\ln(\hat{\sigma}_1^2) + \frac{n_2}{2}\ln(\hat{\sigma}_2^2).
		\end{align}
		Multiplying by $(-2)$ across all terms to obtain the classical test form:
		\begin{align}
			-2 \ln \Lambda &= (n_1 + n_2)\ln(\hat{\sigma}_0^2) - n_1\ln(\hat{\sigma}_1^2) - n_2\ln(\hat{\sigma}_2^2) = n_1\ln(\hat{\sigma}_0^2) - n_1\ln(\hat{\sigma}_1^2) + n_2\ln(\hat{\sigma}_0^2) - n_2\ln(\hat{\sigma}_2^2) \nonumber \\
			&= n_1 \ln\left( \frac{\hat{\sigma}_0^2}{\hat{\sigma}_1^2} \right) + n_2 \ln\left( \frac{\hat{\sigma}_0^2}{\hat{\sigma}_2^2} \right).
		\end{align}
		Wilks' parametric identity is thus formally derived.
		
		\item \textbf{Asymptotic invocation of Wilks' Theorem and comparison with Snedecor's $F$-test:}
		\textbf{Samuel S. Wilks' Theorem} proves that under parametric regularity conditions (where the true null parameter resides in the interior of subspace $\Theta_0$), the statistic $-2\ln\Lambda$ converges in distribution as sample sizes grow infinitely ($n_1, n_2 \to \infty$) to a central Chi-square distribution ($\chi^2_\nu$).
		The degrees of freedom $\nu$ correspond exactly to the \textbf{geometric difference in the dimension of parameter spaces}:
		\begin{align}
			\nu = \dim(\Theta) - \dim(\Theta_0) = 4 - 3 = 1 \text{ degree of freedom.}
		\end{align}
		(In $\Theta$, the four free parameters were $(\mu_1, \mu_2, \sigma_1^2, \sigma_2^2)$, whereas in $\Theta_0$ they were restricted to three: $(\mu_1, \mu_2, \sigma_0^2)$).
		
		\textbf{Conceptual comparison with Snedecor's $F$:} While the Likelihood Ratio Test (LRT) invoked via Wilks is an \textbf{asymptotic} statistic valid only in large samples and extremely sensitive to the normality assumption, \textbf{Snedecor's $F$-test ($F = s_1^2 / s_2^2$) is the exact finite-sample test} that is Uniformly Most Powerful Unbiased (UMPU), making it the gold standard choice for small- or medium-sized experimental designs in analytical testing.
	\end{enumerate}
\end{solucion}

\begin{solucion}
	\begin{enumerate}
		\item \textbf{Proof of the stochastic uniformity of the $p$-value under $H_0$ ($P \sim U(0,1)$):}
		In any upper one-tailed test with a strictly increasing continuous statistic $T$, the observed $p$-value is the upper exceedance probability under the null model: $P(X) = S_0(T(X)) = 1 - F_0(T(X))$.
		We evaluate the cumulative distribution function of this new random variable $P$, denoted by $F_P(x) = P_{H_0}(P \le x)$ for any fixed real number $x \in (0,1)$:
		\begin{align}
			F_P(x) = P_{H_0}\left( S_0(T) \le x \right) = P_{H_0}\left( 1 - F_0(T) \le x \right) = P_{H_0}\left( F_0(T) \ge 1 - x \right).
		\end{align}
		Since the null cumulative distribution function $F_0(t)$ is continuous and monotonically increasing across the real support, it possesses a unique inverse function $F_0^{-1}$. Applying the monotonic inverse operator to both sides of the inner probability inequality without altering inequality direction:
		\begin{align}
			F_P(x) = P_{H_0}\left( T \ge F_0^{-1}(1 - x) \right).
		\end{align}
		By definition of the survival function under the null hypothesis, $P_{H_0}(T \ge t) = 1 - F_0(t)$. Evaluating at $t = F_0^{-1}(1 - x)$:
		\begin{align}
			F_P(x) = 1 - F_0\left( F_0^{-1}(1 - x) \right) = 1 - (1 - x) = x.
		\end{align}
		Thus, we have analytically proven that for every real quantile $x \in (0,1)$:
		\begin{align}
			P_{H_0}(P \le x) = x.
		\end{align}
		This is uniquely and exclusively the cumulative distribution function of a \textbf{Standard Uniform random variable $U(0,1)$}. The theorem irrefutably demonstrates that if the null hypothesis is true, $p$-values spread across the interval with complete flatness ($f_P(p) = 1$). A $p$-value is not the probability of the true parameter ($P(H_0 \mid D)$), but rather a uniform random quantile generated by the testing procedure under null conditions.
		
		\item \textbf{Rigorous concavity and stochastic dominance under $H_a$ ($G_P(u) > u$):}
		When an experiment is conducted under a true alternative hypothesis $H_a$ where the true population parameter exceeds the null point ($\mu_a > \mu_0$), the new distribution of statistic $T$ is governed by $F_a(t) = P_{H_a}(T \le t)$. We evaluate the distribution function of $p$-values in this scenario ($G_P(u) = P_{H_a}(P \le u)$):
		\begin{align}
			G_P(u) = P_{H_a}\left( S_0(T) \le u \right) = P_{H_a}\left( T \ge S_0^{-1}(u) \right) = 1 - F_a\left( S_0^{-1}(u) \right).
		\end{align}
		By the geometric assumption of an alternative with a positive real effect, the density curve of $T$ under $H_a$ is shifted to the right relative to $H_0$. By first-order stochastic dominance for right-shifted variables:
		\begin{align}
			F_a(t) < F_0(t) \quad \forall t \in \mathbb{R} \implies 1 - F_a(t) > 1 - F_0(t) = S_0(t).
		\end{align}
		Evaluating this stochastic dominance inequality precisely at critical point $t = S_0^{-1}(u)$:
		\begin{align}
			G_P(u) = 1 - F_a\left( S_0^{-1}(u) \right) > S_0\left( S_0^{-1}(u) \right) = u \implies G_P(u) > u \quad \forall u \in (0,1).
		\end{align}
		Moreover, differentiating $G_P(u)$ twice using the chain rule with respect to quantile $u$ verifies that the second derivative is strictly negative ($G_P''(u) < 0$), proving that the accumulation curve of $p$-values under $H_a$ is a \textbf{strictly concave-upward function that rises at the origin, stochastically and absolutely dominating the unit diagonal segment of null uniformity}.
		
		\item \textbf{Limiting differential density and monotone likelihood ratio:}
		We differentiate the $p$-value distribution function under $H_a$ ($G_P(u) = 1 - F_a(S_0^{-1}(u))$) using the chain rule to find the true marginal probability density function of the $p$-value, $g_P(u) = G_P'(u)$:
		\begin{align}
			g_P(u) = \frac{d}{du}\left[ 1 - F_a\left( S_0^{-1}(u) \right) \right] = -f_a\left( S_0^{-1}(u) \right) \cdot \frac{d}{du}\left[ S_0^{-1}(u) \right].
		\end{align}
		By the inverse function derivative theorem, knowing that $S_0(t) = 1 - F_0(t) \implies S_0'(t) = -f_0(t)$:
		\begin{align}
			\frac{d}{du}\left[ S_0^{-1}(u) \right] = \frac{1}{S_0'\left( S_0^{-1}(u) \right)} = \frac{1}{-f_0\left( S_0^{-1}(u) \right)}.
		\end{align}
		Substituting into the $p$-value density expression, the negative signs cancel reciprocally:
		\begin{align}
			g_P(u) = -f_a\left( S_0^{-1}(u) \right) \cdot \left( \frac{-1}{f_0\left( S_0^{-1}(u) \right)} \right) = \frac{f_a\left( S_0^{-1}(u) \right)}{f_0\left( S_0^{-1}(u) \right)}.
		\end{align}
		This establishes the exact identity: \textbf{the density in $p$-value space at any point $u$ is identically the Likelihood Ratio evaluated at the corresponding test quantile ($t = S_0^{-1}(u)$)}.
		Evaluating the density at the extreme rejection boundary where the $p$-value approaches zero from the right ($u \to 0^+$):
		\begin{align}
			\lim_{u \to 0^+} g_P(u) = \lim_{t \to \infty} \frac{f_a(t)}{f_0(t)}.
		\end{align}
		In tests with exponential family or normal tail distributions ($f(t) \propto e^{-(t-\mu)^2/2}$), the likelihood ratio $\dfrac{f_a(t)}{f_0(t)}$ grows exponentially to infinity as $t \to \infty$ whenever $\mu_a > \mu_0$.
		Consequently, $\lim_{u \to 0^+} g_P(u) \to +\infty$.
		This explains mathematically why, given any true population discrepancy ($\mu_a > \mu_0$), if sample size $n$ increases sufficiently, the statistical density of the $p$-value collapses into an asymptotic spike of mass concentrated right at $u = 0$, ensuring test statistical power that converges to certainty ($1-\beta \to 100\%$).
	\end{enumerate}
\end{solucion}
